\begin{corollary}[7.9.4]
\label{IV.2.s7.9.4-fr}
Sea $A$ un anillo local noetheriano.
\begin{enumerate}
\item[\emph{(i)}] Si, para todo anillo cociente íntegro $B$ de $A$, se puede
resolver $\operatorname{Spec}(B)$, entonces las fibras formales de $A$ son
regulares.
\item[\emph{(ii)}] Supongamos que, para todo anillo cociente íntegro $B$ de
$A$, y todo anillo íntegro $B'$ que contenga a $B$, que sea una $B$-álgebra finita
y cuyo cuerpo de fracciones sea radicial sobre el de $B$, se pueda
resolver $\operatorname{Spec}(B')$; entonces las fibras formales de $A$ son
geométricamente regulares.
\end{enumerate}
\end{corollary}

\begin{proof}
Como toda fibra formal de $A$ es una fibra formal de un anillo cociente íntegro
de $A$ en el punto genérico de su espectro, está claro que (i) es una consecuencia
inmediata de (7.9.3), y (ii) se deduce de (i) y de (6.7.7).
\end{proof}
